\documentclass[11pt]{article}
\usepackage[T1]{fontenc}
\usepackage{lmodern}
\usepackage{amsmath,amssymb}
\usepackage{booktabs}
\usepackage{array}
\usepackage[a4paper,margin=2.4cm]{geometry}
\usepackage[hidelinks]{hyperref}
\setlength{\parskip}{5pt}
\title{\vspace{-1.6cm}\Large Basics and Definitions\\[2pt]
\large Everything you need before the research plan on waves in a quasicrystal plate}
\author{Dr.\ Vipin \\ Gurugram University}
\date{30 October 2026}
\begin{document}
\maketitle
\vspace{-6mm}

\section{How to use this booklet}

The research plan is short and it assumes a good deal. It assumes you know what a wave is, what a plate does
to a wave, what stress and strain mean, how an electric field behaves inside a crystal, how heat flows when
it is in a hurry, and why a quasicrystal needs a second kind of displacement that no ordinary solid has. None
of that is difficult, but all of it is unfamiliar at the start, and a plan read without it is just a list of
symbols.

So this booklet starts from the beginning and builds up. Sections two to nine are about waves: what they are,
how we count them in time and in space, how fast they go, why a plate turns a continuous spectrum into a
discrete list, and what those modes look like. Sections ten and eleven are about complex numbers, built from
nothing, because the wavenumber in the plan is complex and the word should not be allowed to frighten anyone.
Sections twelve to fourteen are about force in a solid. Sections fifteen to nineteen are about the physics
that belongs to this particular paper: electricity in a crystal, quasicrystals and phasons, heat with a
memory, and the size effect. Section twenty is a table of every symbol in the plan with its unit and a rough
size. Section twenty one is a glossary. Section twenty two tells you how to test your own work.

There are no boundary conditions here, no solution method, and no numbers to compare your answers against.
Those belong to you. This booklet only makes the plan say what it says.

\section{What a wave is}

Strike a long table at one end. The sound reaches the other end almost at once, and yet nothing has travelled
from your hand to the far end except a disturbance. Each small piece of wood pushes the next piece and then
settles back roughly where it was. The pattern of pushing moves; the wood does not.

That is a wave: a disturbance that travels through a material while the material stays at home. Two features
of that sentence matter. First, the wave carries energy, and that is why it can do work at the far end, such
as rattle a cup. Second, the wave's speed is decided by the material and by the geometry, not by how hard you
hit it. Hit the table harder and the sound is louder; its speed and its pitch are unchanged.

A wave has an amplitude, which is the size of the disturbance, and a shape. The shapes we care about are
smooth and repeating, built from sines. This is not laziness. Any smooth repeating shape, however odd, can be
written as a sum of sines of different wavelengths, so a theory that handles one sine handles everything, one
sine at a time. When the plan says that a field varies as a single exponential in space and time, it is
choosing one of those sines and following it.

Solids carry two kinds of disturbance. A push and pull along the direction of travel is a longitudinal wave;
sound in air is of this kind. A sideways sliding across the direction of travel is a transverse wave; light
is of this kind, and so is the shaking that arrives after the first jolt of an earthquake. A solid can carry
both at once, and in a plate the two get mixed together at the free surfaces. That mixing is the origin of
the mode structure we come to later.

\section{Frequency, period, and counting in time}

Frequency is how often the disturbance repeats in one second. The unit is the hertz, one cycle per second. A
guitar string tuned to a low E vibrates about eighty times a second; the high E is some four hundred times a
second. Your ear reads frequency as pitch, which is why the two strings sound different although they are
plucked the same way.

The period is the time for one full cycle, and it is the reciprocal of the frequency. A vibration at four
hundred hertz has a period of two and a half milliseconds. Fast vibration, short period. Slow vibration, long
period. Half a period, by the way, is not a full cycle: a full cycle goes up, down, and back to where it
started, while half a cycle goes up and down once.

Angular frequency is frequency multiplied by two pi, and it is written with the Greek letter omega. The
reason is that one cycle is one full turn, and a full turn is two pi radians. Measuring angles in radians
rather than degrees makes the algebra of oscillations tidy, and the tidiness is worth one extra symbol. When
the plan writes omega, read it as the number of radians of phase per second.

Frequency is not just a label. It sets the rate at which the material is being asked to change, and every
material has a rate beyond which it cannot keep up. A material that follows faithfully at low frequency lags
at high frequency, and a lag is where energy leaks out as heat. Most of what the paper calls attenuation comes
from that lag, in one mechanism or another.

\section{Wavelength, and counting in space}

Wavelength is the distance between two points that are doing the same thing at the same time. It is measured
in metres or fractions of a metre. The wavelength of middle A in air is about three quarters of a metre,
which is why a doorway does not block it; the wavelength of a whistle is a few centimetres, which is why the
sound turns corners reluctantly.

Wavelength matters because a wave only notices obstacles of a size comparable with itself. That is the whole
of diffraction in one sentence, and it applies inside a solid as much as in a room. In a plate, the ratio of
wavelength to thickness decides which kinds of motion are geometrically possible, and that ratio is the true
subject of the paper.

The three quantities are related by one line of arithmetic. A wave travels one wavelength in one period, so
speed is wavelength divided by period, and period is the reciprocal of frequency, so speed is wavelength times
frequency. That is all.

\section{The wavenumber}

The wavenumber counts the repetitions in space the way frequency counts them in time. Its symbol is k and its
unit is one over a metre.

The useful definition is not crests per metre but phase per metre. One full cycle is two pi radians, so a
wave of wavelength lambda has wavenumber two pi divided by lambda. Keeping the two pi makes the travelling
wave read as a cosine of k x minus omega t, with no stray constants anywhere else. Every formula in the paper
is written on that convention, and it is worth adopting it for life.

To feel the difference between the two counts, stand at a railway line. The number of carriages passing per
minute is a frequency. The number of carriages in a hundred metres of track is a wavelength count, and its
reciprocal is a wavenumber. Both describe the same train.

For a plate, the wavenumber along the plate is the important number, because the plate is long in that
direction and thin across it. The wave may pick any wavenumber along the plate, but the plate then decides
which through-thickness shapes are compatible with it. The paper is a careful study of that decision.

\section{Phase velocity, group velocity, and dispersion}

Phase velocity is the speed at which a crest moves: frequency divided by wavenumber, which is the same as
wavelength times frequency. If that were the whole story, waves would be simple.

It is not the whole story because phase velocity usually depends on frequency. That dependence is dispersion,
and it is why a prism spreads white light into a rainbow. In a plate, dispersion is unavoidable, because the
zigzag path between the faces depends on the wavelength.

There is a second speed, the group velocity, which is the speed at which a bundle of waves travels and the
speed at which energy moves. In a dispersive medium the two differ, and a short pulse spreads out as it
travels because its different frequencies travel at different speeds. If you have ever heard a distant clap
arrive as a soft hiss rather than a sharp crack, you have heard dispersion in air.

Dispersion is a nuisance for communication and a gift for measurement. A measured dispersion curve is a
fingerprint of the plate. Change the material, the thickness, or the amount of damage and the fingerprint
moves. That is the basis of a large industry of testing by guided waves, and it is why a paper on a new plate
material is worth writing.

\section{A plate, and why waves stay inside it}

A plate is a solid much thinner in one direction than in the other two. In our problem the thickness is two
h, so the faces are at plus h and minus h, and the plate is taken as infinite in the other directions, which
is the mathematician's way of saying long enough not to matter.

Send a wave into the plate at an angle. It reaches a face; a free face cannot push back into the air in any
useful way, so nearly all of the wave reflects. It crosses to the other face, reflects again, and continues.
The wave is trapped, and while trapped it travels along the plate.

The picture to hold is a zigzag. The wave does not travel straight down the plate; it travels diagonally,
bouncing. Its forward progress is the projection of that diagonal, and the apparent speed along the plate is
therefore less than the speed of the wave in the bulk material. The steeper the zigzag, the slower the
apparent speed.

Two consequences follow immediately, and both shape the paper. First, the thickness appears directly in the
problem, because the round trip length is fixed by it. Second, the two reflections must agree in phase for a
mode to exist at all. That condition is what turns a continuum of possibilities into a discrete list of
modes.

\section{Modes, symmetric and antisymmetric}

A mode is one of the allowed zigzags. Each has a definite shape across the thickness and its own dispersion
relation. You can picture a mode as a standing pattern through the thickness that is carried along the plate
like a wave on a string that is moving sideways at the same time.

The modes come in two families, and symmetry is the reason. The plate looks the same if you turn it over, so
every solution must have a partner under that turn.

In the symmetric family the two faces move in and out together, opposite to one another, and the middle of
the plate moves only along the plate. The lowest such mode stretches and squeezes the plate and is often
called the extensional mode. In the antisymmetric family the two faces move the same way at the same time, so
the plate bends. The lowest of these is the bending mode, and it is what a plate does when you tap it in the
middle.

At low frequency only these two exist. As frequency rises, more appear, each at its own cutoff. In a simple
elastic plate the families do not mix, because a symmetric motion cannot become antisymmetric without
breaking the symmetry of the problem. In our material they can mix at the surfaces, because the electric and
thermal conditions at a face need not respect the same symmetry. Part of the paper's work is to see how much
mixing there is.

\section{Cutoff, and the frequency times thickness axis}

The cutoff of a mode is the lowest frequency at which it can exist. Below it the mode does not travel; it
decays away from where it was created and is called evanescent. The physical reason is the zigzag again: at
low frequency the wavelength is long compared with the thickness, so a wave cannot complete a round trip and
come back in step. Only the lowest two modes escape this restriction, which is why a thin plate driven slowly
carries only two useful modes.

A dispersion curve is a graph of phase velocity against frequency, one curve per mode. Cutoffs appear as the
points where curves begin. The higher modes rise steeply after their cutoff and then flatten towards the
speed of a wave on the surface of a half space, the Rayleigh speed. Any calculation that does not approach
that limit has something wrong with it.

One practical trick removes most of the confusion about thickness. A plate twice as thick at half the
frequency behaves almost identically, because the wave knows the product of frequency and thickness rather
than either alone. So the horizontal axis of every figure is frequency multiplied by thickness, in hertz
metres. On that axis, plates of different thickness and the same material fall on the same curves.

\section{Complex numbers, from nothing}

A complex number is an ordinary number with one extra idea attached, and the idea is a number whose square is
minus one.

Ordinary numbers do not have it: anything real squared is positive or zero. So x squared equals minus one has
no solution among them. We invent one and call it i, with the rule that i squared is minus one. It looks
arbitrary, and for a century it was treated as a trick. It is really a way of writing two numbers at once,
and that is exactly what a wave needs, because a wave has an amplitude and a phase.

A complex number is written a plus b i, where a and b are ordinary numbers. The first is called the real part
and the second the imaginary part. The name imaginary is an accident of history. If a is the length of a step
to the east and b the length of a step to the north, then a plus b i is simply where you end up. In other
words, it is a pair of numbers written compactly.

Addition is easy: add east to east and north to north. Multiplication follows the usual rules with i squared
replaced by minus one. The size, or modulus, of a plus b i is the square root of a squared plus b squared,
which is Pythagoras and the length of the step to the corner. Every complex number can also be written as its
size times a phase, where the phase is the angle of the step from the east direction, measured in radians.

That polar form carries the physics. When you multiply two complex numbers the sizes multiply and the phases
add. A wave travelling through a material has a size that shrinks and a phase that turns, so a single complex
number can carry both. And the exponential of i times an angle equals a cosine plus i times a sine, which
means a wave can be written as a size times an exponential, and then the trigonometry of waves becomes the
arithmetic of exponents.

Two rules follow, and they explain why the equations in the plan look so short. Differentiating with respect
to time multiplies by minus i omega; differentiating twice multiplies by minus omega squared; differentiating
with respect to position multiplies by i k. So a symbol like omega squared or k squared in front of a term is a
derivative in disguise. If you can read that, you can read most of the paper.

\section{The complex wavenumber, and attenuation}

Write a wave travelling forward as a size times the exponential of i times k x minus omega t, and let k be
complex: k equals k real plus i k imaginary. Split the exponential into a phase part and a size part. The
phase part turns, and its rate is set by the real part of k, so the wavelength is two pi divided by the real
part. The size part shrinks with distance, and its rate is set by the imaginary part.

So a complex wavenumber is not a mathematical decoration. Its real part is the wavelength and its imaginary
part is the loss. The convention here is that a positive imaginary part means decay in the positive
direction. That is a choice, like measuring height upwards, but every paper must state it, because a reader
who assumes the opposite will see growth where the author wrote decay.

Attenuation is quoted either as the imaginary part itself, in reciprocal metres, in which case the unit is
called nepers per metre, or in decibels per metre. A neper is like a radian: a ratio wearing the clothes of a
unit. The conversion is about eight and a half decibels to a neper, because decibels use base ten logarithms
and nepers use base e.

The ratio of imaginary to real part is the useful number. It says how much amplitude is lost per radian of
phase, which is the same as the loss per wavelength travelled. A mode with a small ratio goes a long way; a
mode with a ratio near one dies within a wavelength and is useless in a device. When the plan says that an
effect may be too small to see, it usually means that the change in this ratio is smaller than the accuracy of
a measurement.

\section{Stress, strain, and Hooke's law}

Pull a rubber band. It gets longer. The change in length divided by the original length is the strain, a pure
number, usually small. The band pulls back, and the force it transmits divided by the area over which it acts
is the stress, measured in pascals, which are newtons per square metre.

There are two kinds of stress. A normal stress pulls apart or pushes together, and a shear stress slides one
layer over the next. A solid carries both, and in a plate both appear at once, because the zigzag keeps
converting one into the other.

Stress and strain are not single numbers. At a point there are three directions, and a stress can act on each
of three faces, which makes nine components, of which six are independent in a symmetric material. When the
plan writes a symbol with indices like i j, it is naming one component. You do not need to hold nine numbers
in your head; you need to know that the notation is counting them.

Strain can also vary from point to point. The top of a bent ruler is stretched more than the middle. How fast
the strain changes along a direction is the strain gradient, with units of one over length. In ordinary
elasticity the strain gradient costs nothing; in the theory we use it costs a little, and that little is what
makes a thin plate behave differently from a scaled drawing of a thick one.

Hooke's law says that, for small pushes, the response is proportional to the cause. Hang a weight on a spring
and it stretches; double the weight and it stretches twice as much. In a solid the bookkeeping is heavier but
the idea is identical: a strain in one direction produces stresses in several, and the table of proportions is
the stiffness. Pull a strip of rubber and it narrows as it lengthens; the ratio of the narrowing to the
lengthening is Poisson's ratio, near one half for rubber, near zero for cork, which is why cork makes a good
bottle stopper.

For a material that looks the same in all directions the table collapses to two numbers, and it is convenient
to turn those two into two wave speeds: a faster one for push and pull, a slower one for shear. Our materials
are not like that. A crystal has a preferred axis, and a laminated composite is a stack rather than a
material, so its effective stiffness is an average. Averages always involve a choice, and a paper must say
which choice it made.

\section{Balance laws, and what a divergence means}

Physics rests on a few conservation statements. Mass is conserved: what flows into a small box either stays or
leaves. Momentum is conserved: the net force on a box equals the rate of change of its momentum, which for a
solid is its mass times its acceleration. Energy is conserved: what the wave stores as elastic energy and
carries as kinetic energy, it must have taken from somewhere.

These are written with a divergence, a word that alarms people unnecessarily. A divergence is a net outflow.
Imagine a small box and a flow of something through it. If more leaves through the faces than enters, the
amount inside falls and the divergence is positive. If more enters, the divergence is negative and the amount
inside grows. That is all it means.

Put a material law and a balance law together and you get a wave equation. The material law says how stress
follows strain; the balance law says how stress accelerates the material. Combine them and the acceleration
of the displacement turns out to be proportional to its curvature in space. That equation has waves as its
solutions, and their speed is the square root of stiffness over density. Every wave speed in the paper comes
from that one recipe, applied more carefully each time.

\section{Electric fields, potential, and what a crystal does to them}

A charge feels a force when other charges are near. Rather than compute that force charge by charge, we say
the space is filled with an electric field, and a charge placed in the field feels a force equal to its charge
times the field. The field is force per unit charge, and its unit is the volt per metre.

The potential is the second convenience, and it is the one that takes getting used to. It is energy per unit
charge. Moving a charge from one place to another costs energy, and the cost per unit charge is the potential
difference, which we normally call voltage. A battery labelled nine volts promises nine joules for every
coulomb taken through it.

Field and potential are related by a gradient with a minus sign: the field points the way the potential falls
fastest. Water runs downhill; positive charge drifts towards lower potential. The minus sign exists so that
those two sentences agree, and it is the source of half the sign errors in the literature.

Gauss's law says that field lines begin on positive charge and end on negative charge. Draw a closed surface
and count the lines poking out of it, with sign: the count equals the charge inside. Inside a material the
quantity to count is the electric displacement, which is the field scaled by the permittivity, because the
atoms of the material stretch into tiny dipoles and cancel part of the charge you supplied. The displacement
counts only the charge you control, which is why it is the convenient variable in a device.

Our material is piezoelectric as well as being a quasicrystal, so squeezing it produces a voltage and applying
a voltage squeezes it. This coupling is what makes the electric field a full partner in the wave problem
rather than a passenger, and it is why the plan carries the electric potential as one of the unknowns.

\section{What a quasicrystal is}

Most solids are crystals. Their atoms sit on a lattice that repeats, so a small block of the crystal slid by
one lattice spacing looks exactly the same. Quasicrystals do not do that. Their diffraction patterns have
sharp spots, which says they have long-range order, but the spots lie at spacings no repeating lattice can
produce. The first such material was reported in 1982 and the Nobel Prize in Chemistry followed in 2011.

The easiest picture is a tiling. Take two shapes, a fat rhombus and a thin one, and lay them out so that the
pattern never repeats, yet follows the same local rules everywhere. The Penrose tiling of the plane is the
best known example. Stand on any tile: the arrangement is ordered, but no two places are identical in the way
two places in a brick wall are. A quasicrystal is the three-dimensional version of that idea.

Quasicrystals are classified by symmetry, as crystals are. A one-dimensional quasicrystal is periodic in two
directions and aperiodic in the third. If the periodic plane has six-fold symmetry, it is called a
one-dimensional hexagonal quasicrystal, and it is the case with the largest plate literature. In our problem
the aperiodic direction is the thickness, which is a deliberate choice, because the plate faces then become
the place where the unusual field has to be told what to do.

\section{Why a quasicrystal needs two displacements}

In an ordinary solid one field is enough. The displacement tells you how far each point has moved from where
it was, and the strain is its gradient. The waves of that field are sound waves, and the field is called the
phonon displacement.

In a quasicrystal a second thing can happen. A block of tiles can rearrange without any tile being squeezed.
They slide over one another, as a row of books slides along a shelf, and after the slide the material is
locally a different tiling. That rearrangement is described by a second field, the phason displacement.

It is worth being clear about what the phason displacement is not. It is not a movement of matter through
space; no atom packs up and leaves. What changes is the way the tiles fit together, and the phason field
measures how much the fit has changed. That is why the literature speaks of phason strain but not of phason
velocity in the ordinary sense.

The consequence is a richer elasticity. Some constants describe ordinary stretching, some the cost of a
rearrangement, and some the coupling between the two. Once the material is deformed, a phonon strain can
generate a phason strain and the reverse, and that two way traffic is what makes these materials interesting
to a mechanician. In our plate the phason displacement along the thickness is the fourth unknown, and its
coupling to the phonon field is the reason the paper is not a copy of an elastic plate calculation.

\section{Phasons relax, and sometimes they travel}

The phason field has a feature that no ordinary elastic field has: it is not conserved. A rearrangement can
happen and then undo itself. In mechanical language, the phason has friction.

If the friction is large compared with the inertia, a disturbance in the phason field does not travel. It
spreads and dies, exactly as heat spreads through a metal bar. If the friction is small, the phason field
carries a wave with a finite speed. If the two are comparable, the behaviour depends on frequency: slow
disturbances diffuse, fast ones propagate.

All three cases are described by one equation, called the telegraph equation, which contains both an inertia
term and a friction term. Turn the friction off and it becomes a wave equation. Turn the inertia off and it
becomes a diffusion equation. Keep both and you have a crossover, at a frequency set by the ratio of friction
to inertia. This is not a curiosity of our model; it is the standard dynamics of the phason field, and the
literature on quasicrystal hydrodynamics wrote it down in the nineteen eighties.

Why does a plate matter here? Because a plate has a smallest wavelength it can hold, since a half wavelength
cannot be longer than the plate is thick. That sets the lowest frequency at which a phason mode can exist, and
by choosing the thickness one can place the phason mode on either side of the crossover. That is a real
experimental knob, and it is the reason the paper is about plates rather than about bulk material.

\section{Heat, and heat with a memory}

Heat flows from hot to cold, at a rate proportional to the temperature gradient. That is Fourier's law, and
for steady problems it is all you need. Combine it with energy conservation and you get a diffusion equation
for the temperature, in which a point warms if it is colder than its neighbours on average.

The diffusion equation has an uncomfortable feature: it lets a temperature change at one point be felt
everywhere at once, however slightly. For ordinary engineering that is harmless. For a wave at very high
frequency it is not, because the equation then predicts a loss that grows without limit as frequency rises.
Nature does not do that, and the divergence is a warning that the theory is being asked a question outside
its range.

The repair is to give the heat flux a short memory. The flux does not respond instantly to the gradient; it
relaxes towards it over a relaxation time. One memory gives a wave equation for heat. Two memories give the
version introduced by Moore, Gibson and Thompson in 1994 and called the Moore-Gibson-Thompson theory, or MGT.

MGT changes the picture in a way that can be stated without any algebra. At low frequency the relaxation time
multiplied by the frequency is small, and the classical answer comes back. At high frequency the flux lags,
the loss stops rising and levels off at a plateau. That plateau is the observable difference between the two
descriptions, and it is one of the things the paper has to say something about.

Whether the plateau is visible depends on the thickness. The time for heat to diffuse across a plate grows as
the square of the thickness, and the ratio of the relaxation time to that diffusion time tells you whether
the memory matters for that plate. Thin plates are therefore the interesting ones, and very thin plates leave
the range in which a continuum description of heat makes sense at all. That double restriction is why one
figure in the plan is a map of thickness against relaxation time.

\section{The size effect in words}

At very small thickness, a plate stops behaving like a scaled drawing of a thick plate. Two ideas are used to
describe what happens, and both are needed because experiments show both kinds of deviation.

The first is nonlocal elasticity. The stress at a point is allowed to depend on the strain in a neighbourhood
around it. The image is a row of springs linked by weak cross springs: a displacement in one place is felt a
little further along. Nonlocality makes the material softer than the classical theory says.

The second is strain gradient elasticity. The energy is allowed to depend on how quickly the strain changes
from point to point, as well as on the strain itself. The image is a sheet of rubber: a patch of uniform
stretch and a patch of rapidly changing stretch can have the same strain at their centres, but the second
costs more energy. The gradient makes the material stiffer.

The version used here combines the two and comes from Lim, Zhang and Reddy. Each idea brings one length, one
for nonlocality and one for the gradient, and the two compete. Both lengths are written as ratios to the
thickness, because a length on its own tells you nothing and a ratio to the geometry tells you everything.

There is a price, and it is the sentence that closes the model section of the plan. A theory with a second
derivative of the strain raises the order of the equations, and higher order equations need more conditions at
the boundary. A classical plate needs three conditions at each face; this one needs more, and which extra
conditions to use is not settled. That is your decision to make, and the plan says so.

\section{What a model is allowed to assume}

A model is not a photograph of the world. It is a set of statements that are true enough, in a range, to
answer a question. The range is part of the model.

Our model assumes the material is a continuum, so no atoms are followed. It assumes a linear response, so
doubling the cause doubles the effect. It assumes constant properties at a reference temperature, which is
strong, given that temperature is one of the fields. It assumes the phason friction is isotropic in the plane
of the plate, although the literature carries a tensor. It assumes a single phason inertia. It assumes the
faces are free of pressure and of charge.

Each assumption is a place where a referee will push, and each deserves a sentence in the paper. The good news
is that assumptions can be tested. If a conclusion survives when the assumption is relaxed, the assumption was
decorative; if it does not, the reader needs to know.

\section{Table of symbols used in the plan}

\begin{center}
\small
\begin{tabular}{p{1.9cm}p{6.0cm}p{2.4cm}p{2.6cm}}
\toprule
symbol & meaning in words & unit & rough size\\
\midrule
$u_x,u_z$ & phonon displacement along and across the plate & m & $10^{-10}$ to $10^{-8}$\\
$w_z$ & phason displacement along the aperiodic axis & m & $10^{-11}$ to $10^{-9}$\\
$w_{ij}$ & phason strain, gradient of the phason displacement & none & $10^{-6}$ to $10^{-4}$\\
$\varepsilon_{ij}$ & phonon strain & none & $10^{-6}$ to $10^{-4}$\\
$\sigma_{ij},t_{ij}$ & local stress, total nonlocal stress & Pa & $10^{5}$ to $10^{8}$\\
$H_{ij}$ & phason stress & Pa & $10^{4}$ to $10^{7}$\\
$C_{ijkl}$ & phonon elastic stiffness & Pa & $10^{10}$ to $10^{11}$\\
$K_{ijkl}$ & phason elastic stiffness & Pa & $10^{9}$ to $10^{10}$\\
$R_{ijkl}$ & phonon-phason coupling & Pa & $10^{9}$ to $10^{10}$\\
$\rho$ & mass density & kg\,m$^{-3}$ & $4\times10^{3}$ to $6\times10^{3}$\\
$\rho_p$ & phason mass density & kg\,m$^{-3}$ & $10^{-3}$ to $10^{-1}$ times $\rho$\\
$\gamma$ & phason friction coefficient & kg\,m$^{-3}$s$^{-1}$ & $10^{5}$ to $10^{9}$\\
$\varphi$ & electric potential & V & 0.1 to 10\\
$E_i$ & electric field, minus the gradient of the potential & V\,m$^{-1}$ & $10^{4}$ to $10^{7}$\\
$D_i$ & electric displacement & C\,m$^{-2}$ & $10^{-3}$ to $10^{-1}$\\
$e_{kij}$ & piezoelectric constant & C\,m$^{-2}$ & 0.1 to 2\\
$\epsilon_{ik}$ & permittivity & F\,m$^{-1}$ & $10^{-10}$ to $10^{-9}$\\
$p_i$ & pyroelectric coefficient & C\,m$^{-2}$K$^{-1}$ & $10^{-6}$ to $10^{-4}$\\
$\theta$ & temperature rise above the reference value & K & below 1\\
$T_0$ & reference temperature & K & about 300\\
$\beta_{ij}$ & thermal stress, stiffness times thermal expansion & Pa\,K$^{-1}$ & $10^{4}$ to $10^{6}$\\
$\beta^{w}_{ij}$ & thermo-phason coupling, no published source & Pa\,K$^{-1}$ & unknown\\
$\alpha_{kl}$ & thermal expansion coefficient & K$^{-1}$ & $10^{-6}$ to $10^{-5}$\\
$\rho c_p$ & heat capacity per unit volume & J\,m$^{-3}$K$^{-1}$ & $10^{6}$ to $3\times10^{6}$\\
$k_{\rm th}$ & thermal conductivity & W\,m$^{-1}$K$^{-1}$ & 1 to 10\\
$\kappa$ & thermal diffusivity & m$^2$s$^{-1}$ & $10^{-7}$ to $10^{-6}$\\
$\tau_q$ & relaxation time of the heat flux (MGT) & s & $10^{-12}$ to $10^{-10}$\\
$\tau_T$ & diffusion time across the plate & s & $10^{-13}$ to $10^{-12}$\\
$l_{\rm nl}$, $l_g$ & nonlocal and gradient lengths & m & $10^{-9}$ to $10^{-7}$\\
$2h$, $h$ & plate thickness and half of it & m & $10^{-9}$ to $10^{-3}$\\
$k$ & complex wavenumber along the plate & m$^{-1}$ & $10^{4}$ to $10^{8}$\\
$k_r,k_i$ & real part sets wavelength, imaginary part sets decay & m$^{-1}$ & $k_i\ll k_r$\\
$\lambda$ & wavelength & m & $10^{-8}$ to $10^{-4}$\\
$\omega$ & angular frequency & rad\,s$^{-1}$ & $10^{6}$ to $10^{11}$\\
$f$ & ordinary frequency & Hz & $10^{5}$ to $10^{10}$\\
$f(2h)$ & frequency times thickness, the comparison axis & Hz\,m & 1 to $10^{4}$\\
\bottomrule
\end{tabular}
\end{center}

\section{Glossary}

\textit{Aperiodic} means ordered but not repeating.

\textit{Coupling} is any term through which one field influences another.

\textit{Cutoff} is the lowest frequency at which a mode exists.

\textit{Dispersion} is the dependence of phase velocity on frequency.

\textit{Dissipation} is the irreversible conversion of wave energy into heat. A positive friction coefficient
is what makes dissipation positive.

\textit{Phason} is the field describing local rearrangements of the tiling. \textit{Phonon} is the ordinary
displacement field.

\textit{Phason strain} is the gradient of the phason displacement. \textit{Phason stress} is its conjugate.

\textit{Piezoelectricity} is the coupling between strain and electric polarisation.

\textit{Relaxation time} is how long a system takes to come back to equilibrium after a small push.

\textit{Strain} is change of length per unit length; \textit{stress} is force per unit area.

\textit{Telegraph equation} is the equation containing both inertia and friction, with diffusion at low
frequency and waves at high frequency.

\textit{Unit cell} is the repeating block of a crystal. A quasicrystal has none, which is the whole point.

\section{How to check your own work}

Five generic tests will catch most mistakes, and none of them requires a published number.

\textit{First, units.} Write the unit of every term in every equation before coding. Terms that disagree are
usually wrong in sign as well. Keep the table and put it in the paper.

\textit{Second, limits.} Every new term has a limit in which it vanishes. Put the phason friction to zero and
you must recover a lossless phason wave. Put the phason inertia to zero and you must recover diffusion. Put
the thermal coupling to zero and the attenuation must fall to its mechanical value. Compute the reduced case
separately and compare, do not simply delete a line from your own code.

\textit{Third, symmetry.} Turn the plate over. The dispersion curves must be unchanged for each family. If a
mode that should be symmetric comes out mixed, either the coupling is real or the assembly has a sign error,
and separating those two possibilities is worth a day of work.

\textit{Fourth, switch a term off in two ways.} Set a coupling constant to zero, and separately assemble a
model in which the term was never written. The two answers must agree to six decimal places. This tests your
code, not the physics.

\textit{Fifth, compare with a cited source.} Digitise two or three points from a published curve and compare
shape and order. The constants will differ, so you are looking for the same trend, not the same digits. If
your curve bends the other way, stop and find out why before writing a sentence.

Three habits complete the set. Keep a log of every run so that results can be reproduced. Plot whatever you
are about to quote, because a table hides a shape. And repeat any surprising result by a different route
before believing it.

\end{document}
